\section{検定の数理 : 統計量と自由度はどこから出てくるか}

    検定の章では, $t$統計量がおよそ2.25, $p$値がおよそ0.028という数字を示しただけで, どう求めたのかは伏せてあった.
    カイ二乗統計量が13.64, $p$値がおよそ0.0035と出たところも同じである.
    その計算と, 比べる先の分布がどこから来るのかを, ここにまとめてある.

  \subsection{自由度は何を数えているのか}

    ここで扱う検定には, \textbf{自由度}という量が繰り返し登場する.
    $t$分布の形も, カイ二乗分布の形も, 自由度で決まる.
    つまり, 自由度を取り違えると, 同じ検定統計量からでも違う$p$値が出てしまう.
    $p$値に直結する量でありながら, 自由度が何を数えているのかは, 式の形からは読み取れない.
    その意味が最もはっきり見えるのは, 母集団の分散を標本から推定する場面である.

    \subsubsection*{母集団の分散を推定したい. ところが母平均がわからない}

    検定のあちこちで必要になるのが, 母集団の分散（母分散）$\sigma^2$ の推定だ.
    母分散は, 母集団の各データが母平均 $\mu$ からどれだけずれているかを二乗して, 母集団全体で平均した量である.
    では, 手元の標本からこの量を推定したいのだが, ここで困ったことになる.
    定義の中身に登場する $\mu$ がわからない.
    そもそも母集団のことが直接わからないからこそ, 標本を使って推定しようとしているのだった.
    それなのに, 推定したい量の定義の中に, もう一つの「わからない」量が入っている.

    ならば, $\mu$ の代わりに, 手元の標本から計算できる量を使う.
    最も自然な代用品は, 標本平均 $\bar{X} = \frac{1}{n}\sum X_i$ だ.
    そこで $\mu$ を $\bar{X}$ で置き換えた量
    \[
      \frac{1}{n} \sum_{i=1}^n (X_i - \bar{X})^2
    \]
    を考えてみたくなる.
    データを要約するときに定義した分散と同じ形だし, これでうまくいきそうな気もする.
    ただ, $\bar{X}$ はあくまで $\mu$ の代用品だから, 置き換えたことで何かがずれていないかは確かめておきたい.

    \subsubsection*{\texorpdfstring{$\bar{X}$}{X̄} で置き換えると, 何がずれるのか}

    母分散の分子に当たる $\sum(X_i - \mu)^2$ と, 標本から計算する $\sum(X_i - \bar{X})^2$ の差を調べる.
    $X_i - \mu = (X_i - \bar{X}) + (\bar{X} - \mu)$ と分解してから二乗を展開すると,
    \begin{align*}
      \sum_{i=1}^n (X_i - \mu)^2
      &= \sum_{i=1}^n \left[ (X_i - \bar{X}) + (\bar{X} - \mu) \right]^2 \\
      &= \sum_{i=1}^n (X_i - \bar{X})^2 + 2(\bar{X} - \mu) \sum_{i=1}^n (X_i - \bar{X}) + \sum_{i=1}^n (\bar{X} - \mu)^2
    \end{align*}
    となる.
    第二項の和の部分 $\sum (X_i - \bar{X})$ は $\bar{X}$ の定義から常にゼロになるので, 第二項そのものが消えて
    \[
      \sum_{i=1}^n (X_i - \mu)^2 = \sum_{i=1}^n (X_i - \bar{X})^2 + n(\bar{X} - \mu)^2
    \]
    が得られる.

    左辺は, 各データが\textbf{母平均から}どれだけずれているかの総量で, 本当に知りたい量である.
    第一項は, 各データが\textbf{標本平均から}どれだけずれているかの総量で, 手元で計算できる.
    残る第二項は, 代用品の $\bar{X}$ 自身が本物の $\mu$ からどれだけずれているかを, データ全体分（$n$倍）まとめた量だ.
    どのくらいの大きさなのか, 期待値を取って測ってみよう.

    \subsubsection*{ずれの大きさはデータ何個分にあたるか}

    左辺の期待値は, 母分散の定義から各 $X_i$ について $E[(X_i - \mu)^2] = \sigma^2$ なので
    \[
      E\!\left[\sum_{i=1}^n (X_i - \mu)^2\right] = n\sigma^2
    \]
    である.

    第二項はどうか.
    標本平均 $\bar{X}$ 自体が母平均からどれだけずれるかは標準誤差で測れて, その分散は $\sigma^2/n$ だった.
    つまり $E[(\bar{X} - \mu)^2] = \sigma^2/n$ となるから,
    \[
      E[n(\bar{X} - \mu)^2] = n \cdot \frac{\sigma^2}{n} = \sigma^2
    \]
    である.

    第二項の期待値は, データひとつ分の分散 $\sigma^2$ ちょうどである.
    第二項は, $\mu$ の代わりに使った $\bar{X}$ 自身のずれから生じた項だった.
    平均の推定と分散の推定を同じ $n$個のデータで済ませようとした代償が, 期待値で測ると, データひとつ分の分散にあたる.

    \subsubsection*{自由度 \texorpdfstring{$n-1$}{n-1} の意味}

    左辺の期待値が $n\sigma^2$, 第二項の期待値が $\sigma^2$ なので, 第一項の期待値は引き算で
    \[
      E\!\left[\sum_{i=1}^n (X_i - \bar{X})^2\right] = n\sigma^2 - \sigma^2 = (n-1)\sigma^2
    \]
    と求まる.

    手元で計算できる $\sum(X_i - \bar{X})^2$ が持っているずれの情報は, 期待値でみると, 本当に欲しかった $n\sigma^2$ よりデータひとつ分だけ少ない.
    だから, この量を $n$ で割ると, 母分散より少し小さめの値が出やすい.
    一方, $n-1$ で割れば, 期待値はちょうど $\sigma^2$ に一致する.
    推論でずっと使ってきた不偏分散
    \[
      u^2 = \frac{1}{n-1} \sum_{i=1}^n (X_i - \bar{X})^2
    \]
    が $n$ ではなく $n-1$ で割る形をしているのは, このためである.

    手元の $n$個のデータから出発して, $\bar{X}$ という一つの代表値に集約する過程で, 使える情報がデータひとつ分だけ目減りした.
    残った情報の個数が, データの数 $n$ から一つ引いた $n-1$ である.
    この数を\textbf{自由度}と呼ぶ.

    極端な場合で確かめておこう.
    データが1個しかないとき, 標本平均はそのデータ自身になり, ずれの総量 $\sum(X_i - \bar{X})^2$ は必ずゼロになる.
    ゼロを $n = 1$ で割れば「分散はゼロ」, つまり「母集団はまったくばらついていない」という見積もりが出てくるが, データ1個でばらつきの話ができるはずがない.
    一方, $n - 1 = 0$ で割ろうとすると, 値は定まらない.
    ばらつきの情報が何も残っていないから見積もれない, という答えであり, 「分散はゼロ」よりも実態に合っている.
    平均の推定に1個分の情報を使い切ると, ばらつきに回せる情報はゼロになる. 自由度ゼロとは, そういう状態のことだ.

    \subsubsection*{標本平均を固定して, 残りを動かしてみる}

    自由度 $n-1$ は, 別の角度からも数えられる.
    標本から $\bar{X}$ を計算してその値を一旦固定し, $\frac{1}{n}\sum X_i = \bar{X}$ という条件を満たしたまま $X_1, \ldots, X_n$ を動かしてみる.
    この条件のもとでは, $X_1, \ldots, X_{n-1}$ までは自由に動かせる.
    しかし最後の $X_n$ は
    \[
      X_n = n\bar{X} - \sum_{i=1}^{n-1} X_i
    \]
    と自動的に決まってしまう.
    平均を $\bar{X}$ に合わせるためには, 最後の一つで帳尻を合わせるしかないからだ.
    自由に動かせる個数は $n-1$ にとどまる.

    元のデータは $n$個とも自由に観測されるのだから, この $n-1$ は, データそのものに最初から制約があるという話ではない.
    そこから $\bar{X}$ という代表値を計算して固定したうえで, 残ったずれ $(X_i - \bar{X})$ を眺めると, ずれの組が自由に動ける余地は $n-1$個分しかない.
    「自由度」という名前は, ここから来ている.

    集約の代償として期待値に現れたデータひとつ分の分散と, ずれが自由に動ける個数が一つ少ないことは, 同じ事情の別の表れである.
    標本から計算した値に合わせるという条件が一つ増えるたびに, 自由に残る情報は一つ減る.
    同じ勘定が, このあとの$t$検定にもカイ二乗検定にも出てくる.


  \subsection{\texorpdfstring{$t$}{t}統計量と\texorpdfstring{$p$}{p}値を計算する}

    原作の有無で平均視聴回数に差があるかを$t$検定で調べたとき, 出てきたのは$t$統計量と$p$値の数字だけだった.
    使ったのは, 差の信頼区間を作ったときと同じ2群のデータである.
    \begin{itemize}
      \item \textbf{原作あり群}（$n_1 = 37$本）：平均視聴回数 $= 274.1$万回, 標準偏差 $u_1 = 260.637$万回
      \item \textbf{原作なし群}（$n_2 = 43$本）：平均視聴回数 $= 153.7$万回, 標準偏差 $u_2 = 209.680$万回
    \end{itemize}
    このデータから, 「原作の有無で平均視聴回数に差がある」と言えるかどうかを, どちらが大きいかは問わない両側検定で調べる.

    \subsubsection*{\texorpdfstring{$t$}{t}統計量の定義と計算}

    $t$検定では, 二つの標本平均の差が, ばらつきに比べてどれだけ大きいかを測る.
    各群のばらつきが違いうることを認める立場で, 次の統計量を用いる：
    \[
      t = \frac{\bar{x}_1 - \bar{x}_2}{\sqrt{\dfrac{u_1^2}{n_1} + \dfrac{u_2^2}{n_2}}}
    \]
    分母は, 二つの平均の差の信頼区間を作ったときに登場した式と同じで, 各群の分散をそれぞれの標本サイズで割って足し合わせ, 平方根をとったものだ.
    分子の「差そのもの」を, 分母の「差のばらつきの大きさ」で割ることで, ばらつきを基準に差の大きさを測り直している.
    $u_1$ と $u_2$ を最後まで別々に扱うこの分母が, 両群の母分散が等しいとは仮定しないWelchの$t$検定の分母にあたる.

    数値を代入して計算してみよう.
    平均視聴回数は, 小数第3位までとれば $\bar{x}_1 = 274.082$万回, $\bar{x}_2 = 153.717$万回である.
    まず, 分母の中身を求める：
    \[
      \frac{u_1^2}{n_1} + \frac{u_2^2}{n_2} \approx 1835.984 + 1022.458 = 2858.442
    \]
    その平方根が分母になる：
    \[
      \sqrt{2858.442} \approx 53.464
    \]
    したがって$t$統計量は
    \[
      t = \frac{274.082 - 153.717}{53.464} \approx \frac{120.365}{53.464} \approx 2.251
    \]
    となる.

    \subsubsection*{自由度と\texorpdfstring{$p$}{p}値}

    $t$分布を使って$p$値を求めるには, 自由度を決める必要がある.
    Welchの$t$検定では, 自由度を次の式で計算する：
    \[
      \nu = \frac{\left( \dfrac{u_1^2}{n_1} + \dfrac{u_2^2}{n_2} \right)^2}
                  {\dfrac{(u_1^2/n_1)^2}{n_1 - 1} + \dfrac{(u_2^2/n_2)^2}{n_2 - 1}}
    \]
    ずいぶん込み入った式である.
    なぜこんな式になるのか, そして「自由度」と呼んでいるのに整数ですらない値が出てくるのはどういうことなのか.
    その事情はあとの節で扱うことにして, ここでは, この式に従って計算するとどうなるかだけ確認しておく.

    分子：
    \[
      (1835.984 + 1022.458)^2 = 2858.442^2 \approx 8{,}170{,}691
    \]
    分母：
    \[
      \frac{1835.984^2}{36} + \frac{1022.458^2}{42}
      \approx 93{,}634 + 24{,}891 = 118{,}525
    \]
    したがって
    \[
      \nu \approx \frac{8{,}170{,}691}{118{,}525} \approx 68.936
    \]
    である.

    この$t$統計量 $t \approx 2.251$ を, 自由度およそ$68.936$の$t$分布と比べる.
    両側検定では原作あり群が高い場合も原作なし群が高い場合もまとめて数えるから, $p$値は $|t| \geq 2.251$ となる確率で,
    \[
      p = P(|T_{68.936}| \geq 2.251) \approx 0.0276
    \]
    である.
    有意水準を $\alpha = 0.05$ と決めていたとすると, $0.0276 < 0.05$ なので帰無仮説は棄却される.
    「原作の有無で平均視聴回数に差がある」という結論は, この計算から出ていた.

    差の95\%信頼区間がゼロを含むかどうかと, 両側検定の結果とは対応していた.
    信頼区間を作ったときは倍率を「およそ2」と概算したが, いまは倍率を正確に決められる.
    自由度68.936の$t$分布で真ん中95\%を切り取る倍率は約1.995なので, 差の95\%信頼区間は
    \[
      120.365 \pm 1.995 \times 53.464
    \]
    すなわち, およそ13.7万回から227.0万回となる.
    この区間はゼロを含まない.
    そして, 両側検定の$p$値 $0.0276$ は $0.05$ を下回った.
    どちらも同じ標準誤差と同じ$t$分布から計算しているから, この二つの判定は必ず一致する.

    \subsubsection*{なぜ正規分布ではなく\texorpdfstring{$t$}{t}分布なのか}

    ここで比べた先が正規分布ではなく$t$分布なのは, 分母に標本から計算した分散が入っているからだ.
    もし母分散がわかっていれば, 分子だけを標準化して正規分布で扱えばよかった.
    ところが母分散は未知で, 手元の不偏分散で代用するしかない.
    分母が偶然小さく出た標本では, 同じ差でも割った結果が大きく膨らむ.
    その膨らみのぶん, $t$の分布は正規分布より裾が重くなる.
    その重さを正面から扱った分布が$t$分布だ.

    自由度が大きくなると, 不偏分散のばらつきが小さくなり, $t$分布は正規分布に近づいていく.
    だから標本が大きいときは, $t$分布で計算しても正規分布で計算しても, $p$値はほとんど変わらない.
    逆に標本が小さいときほど, 正規分布で済ませると, 裾の重さを無視して「実際より珍しい」と判定してしまう.


  \subsection{カイ二乗統計量と\texorpdfstring{$p$}{p}値を計算する}

    カテゴリデータの検定では, 血液型と性格に関係があるのかをカイ二乗検定で調べた.
    100人に血液型と性格タイプを尋ねた結果は, 次の通りだった：

    \begin{center}
      \begin{tabular}{l|cccc|c}
            & A型 & B型 & O型 & AB型 & 合計 \\
      \hline
      几帳面 & 18  & 6   & 10  & 8    & 42 \\
      大雑把 & 7   & 19  & 15  & 17   & 58 \\
      \hline
      合計   & 25  & 25  & 25  & 25   & 100
      \end{tabular}
    \end{center}

    帰無仮説「血液型と性格タイプは無関係である」のもとで各マスに入るはずの人数, つまり期待度数は
    \[
    \text{期待度数} = \frac{\text{行の合計} \times \text{列の合計}}{\text{全体の合計}}
    \]
    で求められるのだった.
    「A型で几帳面」なら $42 \times 25 / 100 = 10.5$ である.
    どの血液型も25人ずつだから, 同様に計算すると, 几帳面の行はどのマスも10.5, 大雑把の行はどのマスも14.5 になる.

    \subsubsection*{カイ二乗統計量の計算}

    カイ二乗統計量 $\chi^2$ は, 観測度数と期待度数のずれの大きさを, 次の式で集計する：
    \[
    \chi^2 = \sum_{\text{すべてのマス}} \frac{(O - E)^2}{E}
    \]
    ここで $O$ は観測度数, $E$ は期待度数である.
    ずれを二乗するのは, 正負のずれが打ち消しあってしまうのを防ぐためだ.
    ずれの二乗 $(O - E)^2$ をそのまま足さずに $E$ で割っているのは, マスごとの規模の違いを揃えるためである.
    期待度数が100人のマスで5人ずれるのと10人のマスで5人ずれるのとでは珍しさがまるで違い, 見込みが10人しかないほうがずっと珍しい.

    各マスについてこの値を計算してみよう（小数第3位で四捨五入）：
    \begin{itemize}
      \item $(18 - 10.5)^2 / 10.5 = 56.25 / 10.5 \approx 5.36$
      \item $(6 - 10.5)^2 / 10.5 = 20.25 / 10.5 \approx 1.93$
      \item $(10 - 10.5)^2 / 10.5 = 0.25 / 10.5 \approx 0.02$
      \item $(8 - 10.5)^2 / 10.5 = 6.25 / 10.5 \approx 0.60$
      \item $(7 - 14.5)^2 / 14.5 = 56.25 / 14.5 \approx 3.88$
      \item $(19 - 14.5)^2 / 14.5 = 20.25 / 14.5 \approx 1.40$
      \item $(15 - 14.5)^2 / 14.5 = 0.25 / 14.5 \approx 0.02$
      \item $(17 - 14.5)^2 / 14.5 = 6.25 / 14.5 \approx 0.43$
    \end{itemize}
    これらをすべて足し合わせると,
    \[
    \chi^2 \approx 5.36 + 1.93 + 0.02 + 0.60 + 3.88 + 1.40 + 0.02 + 0.43 = 13.64
    \]
    となる.

    \subsubsection*{自由度と\texorpdfstring{$p$}{p}値}

    この統計量を比べる先はカイ二乗分布で, その形も$t$分布と同じく自由度で決まる.
    クロス集計表の検定では, 自由度は
    \[
    (\text{行数} - 1) \times (\text{列数} - 1) = (2 - 1) \times (4 - 1) = 3
    \]
    である.

    なぜ行数や列数から1を引くのか.
    期待度数を作るときに行の合計と列の合計を標本から計算して使ったので, 表のマスを自由に埋められる余地が減っている.
    実際, この2行4列の表では, たとえば「几帳面」の行のA型・B型・O型の3マスを決めると, 残りのマスはすべて行と列の合計から自動的に決まってしまう.
    自由に決められるマスの数は $(2-1) \times (4-1) = 3$ 個しかない.
    平均を固定すると最後の一つは帳尻合わせで決まる, という勘定と同じことが, ここでは行と列の両方向で起きている.

    したがって, $\chi^2 = 13.64$ を比べる先は, 自由度3のカイ二乗分布になる.
    この分布で13.64以上の値が出る確率を求めると,
    \[
    p = P(\chi^2_3 \geq 13.64) \approx 0.0035
    \]
    である.
    有意水準を $0.05$ としていたとすると $p$値はそれよりはるかに小さく, 帰無仮説は棄却される.
    「血液型と性格タイプは無関係とは言いがたい」という結論も, この計算から出ていた.

    \subsubsection*{なぜカイ二乗分布と比べられるのか}

    カイ二乗分布は, 理論的には, 平均0, 分散1の正規分布に従う独立な変数を二乗して足し合わせたものが従う分布として定義され, 足し合わせる個数が自由度になる.
    観測度数と期待度数のずれは, 標本が大きければ近似的に正規分布に従う.
    カイ二乗統計量は, そのずれを期待度数の大きさで揃えてから二乗して足したもので, 自由に動ける個数分, つまり自由度3のカイ二乗分布で近似できる.
    この導出の細部は本書の範囲を超えるが, 「統計量を, 帰無仮説のもとでの分布と比べて珍しさを測る」という枠組み自体は, $t$検定とまったく同じである.


  \subsection{Welchの自由度はなぜ整数でないのか}

    $t$統計量の$p$値を求めるときに使った自由度 $\nu \approx 68.936$ には, 端数がついていた.
    自由度は, 平均を決めたぶんだけ減った情報の個数だったから, 整数で出てくるはずの量である.
    個数のはずのものに, なぜ端数がつくのか.

    \subsubsection*{ゴセットが置いた仮定}

    古典的な$t$検定を考案したのは, 検定の章のコラムに登場したゴセットである.
    分母に手元の分散を使うと値の分布は正規分布より裾が重くなる. その事情を正面から扱うために考え出されたのが$t$分布だった.
    ただし, この構成が成り立つように, ゴセットは一つ仮定を置いた.
    \textbf{両群の母分散が等しい}という仮定である.
    この仮定があれば, 二つの群のばらつきを一つにまとめて扱えて, 自由度 $n_1 + n_2 - 2$ という整数値をもつ$t$分布が出てくる.
    「使える情報の個数」という解釈が, ここではきれいに整数のまま通る.

    \subsubsection*{仮定を外すと, \texorpdfstring{$t$}{t}分布の足場が崩れる}

    本書が採るWelchの$t$検定は, この等分散の仮定を外したものだった.
    ゴセットの分母は, 等分散の仮定を頼りに二つの群のばらつきを一つの推定量にまとめてあり, 一つにまとまっているからこそ, 全体が自由度 $n_1 + n_2 - 2$ のカイ二乗分布の枠にきれいに収まる.
    一方, Welchの分母の中身 $u_1^2/n_1 + u_2^2/n_2$ は, 二つの独立な不偏分散を, 別々のまま, それぞれ違う重みで足し合わせたものだ.
    重みの違うものどうしを足した和は, 一つのカイ二乗分布にはまとまらない.
    そして, 比である $t$ も, 厳密には$t$分布に従わなくなる.
    仮定を一つ外しただけのつもりが, $t$分布で扱うための足場が一気に崩れてしまうのだ.

    \subsubsection*{期待値と分散を合わせて寄せる}

    まったく新しい分布を作り直すことも考えられるが, それでは新しい分布の性質を一から調べ直すことになる.
    厳密には$t$分布でないとしても, 近似的に$t$分布として扱えれば実用上はそれで足りる.
    では, 何を手がかりに寄せればよいのか.

    手がかりは, 正規分布にある.
    正規分布は, 期待値と分散の二つだけで形が完全に決まる, 際立った分布だった.
    ふつうの分布なら期待値と分散が同じでも形はいくらでも違いうるのに, 正規分布だけはそうならない.
    この経験を引きずって考えれば, 「分布どうしを寄せたいなら, とりあえず期待値と分散を合わせてみる」というのは, それなりに筋が通って見える.

    実際, この見当は当たる.
    Welchの分母の中身, つまり重みの違う二つの不偏分散の和を, 形のよくわかっている「定数倍したカイ二乗分布」で置き換える.
    置き換え先で決めるのは, 定数倍の大きさと自由度の二つである.
    条件のほうも, 置き換え先の期待値と分散を元の和のそれぞれに一致させるという二つで, 決めるものと同じ数だから, ちょうど解ける.
    この連立を解くと, 自由度は母分散を使った式として求まる.
    その母分散を手元の不偏分散で置き換えたのが, 端数のついた $\nu \approx 68.936$ を出した式である.
    この近似は, 考案者の名前からSatterthwaiteの近似と呼ばれる.
    端数のついた値を$t$分布に入れてよいのかが引っかかるかもしれないが, $t$分布の形を定める式の中で $\nu$ は階乗を実数へ滑らかに延ばした関数（ガンマ関数）を介して使われているので, 正の実数なら何を入れても確率分布として一貫した意味がある.
    実際, $t \approx 2.251$ に対する両側の$p$値は, 自由度68でも68.936でも69でも小数第3位までは同じ0.028で, 小数第5位まで見てはじめて0.02760, 0.02756, 0.02755と分かれる.
    丸めても結論は動かないが, 丸めなければならない理由もない.

    もっとも, ここまで「厳密には」と言ってきたのは, 各群の母集団がそれぞれ正規分布に従う場合の話である.
    その導出では, 分子の $\bar{x}_1 - \bar{x}_2$ が正規分布に従うことを使っている.
    実際の母集団がきれいな正規分布であることは期待できないが, 標本サイズがそれなりに大きければ, 中心極限定理によって各群の標本平均は正規分布に近づき, その差もまた正規分布に近づく.
    一方で, 標本サイズが小さく, かつ元の母集団が大きく歪んでいる場合や, 外れ値が平均を強く動かしている場合には, この近似が効かなくなり, Welchの結果も信頼しにくくなる.

    \subsubsection*{出てきた式を眺めてみる}

    各群のばらつきの寄与 $w_i = u_i^2/n_i$ を主役にして, 自由度の式を書き直してみると,
    \[
      \nu = \frac{(w_1 + w_2)^2}{\dfrac{w_1^2}{n_1 - 1} + \dfrac{w_2^2}{n_2 - 1}}
    \]
    という形になる.
    分母では, 各群の寄与の二乗を, その群の自由度で割ってから足している.
    だから片方の群の寄与が飛び抜けて大きければ, $\nu$ はほぼその群の自由度になる.
    両群のばらつきの寄与が等しい, つまり $w_1 = w_2$ とおくと, 分子は $4w_1^2$, 分母は $w_1^2\left(\dfrac{1}{n_1-1} + \dfrac{1}{n_2-1}\right)$ となって $w_1$ が約分で消え, $\nu$ は標本サイズだけで決まる.
    さらに $n_1 = n_2$ なら, ちょうど $n_1 + n_2 - 2$ となり, 古典的な$t$検定の自由度と一致する.
    ばらつきの寄与も標本サイズも同じという, ゴセットの想定に最も近い場合には, ゴセットの自由度がそのまま現れるのだ.

    一方, ばらつきの寄与に偏りがあったり標本サイズが違ったりすると, $\nu$ は整数にならない.
    重みも標本サイズの比も実数値だから, それらを混ぜ合わせた結果が整数になる必然性はない.
    端数が出るのが普通なのだ.

    \subsubsection*{「自由度」という言葉が二つの意味を持っている}

    ここまでを並べると, 「自由度」という言葉が二つの意味で使われていることがわかる.
    はじめに見た自由度は, 平均を決めたぶんだけ減った情報の個数で, 出てくる数は必ず整数だった.
    Welchの自由度は, $t$分布族の中で近似先の位置を指定する目盛りの読みであって, 代償でも個数でもない.
    「自由度が68.936」は, 使える情報が68.936個分ある, という意味ではない. 言葉が同じなだけで, 指しているものが違う.
